% Generated draft: requires mathematical and source review.
\documentclass[11pt]{article}
\usepackage[margin=25mm]{geometry}
\usepackage{fontspec}
\setmainfont{Latin Modern Roman}
\usepackage{amsmath,amssymb,mathtools,unicode-math}
\setmathfont{Latin Modern Math}
\usepackage{enumitem,needspace,xurl,hyperref}
\setlength{\parindent}{0pt}
\setlength{\parskip}{5pt}
\setlength{\emergencystretch}{3em}
\newcommand{\N}{\mathbb{N}}
\newcommand{\Z}{\mathbb{Z}}
\newcommand{\Q}{\mathbb{Q}}
\newcommand{\R}{\mathbb{R}}
\newcommand{\C}{\mathbb{C}}
\newcommand{\F}{\mathbb{F}}
\newcommand{\A}{\mathbb{A}}
\newcommand{\PP}{\mathbb{P}}
\newcommand{\E}{\mathbb{E}}
\newcommand{\eps}{\varepsilon}
\newcommand{\dd}{\,\mathrm{d}}
\begin{document}
\section*{UnsolvedMath ID: 1900019}
\textbf{Source problem code:} catalogue entry
\subsection{Definitions and mathematical statement}
In the upper half-plane model of $\mathbb H^2$, the Farey tessellation is the ideal-triangle tessellation with vertices $\Q\cup\{\infty\}$ at infinity. Reduced fractions $p/q$ and $r/s$, with $\infty=1/0$, are joined by a hyperbolic geodesic precisely when
\[
|ps-qr|=1.
\]
The tessellation is invariant under $\operatorname{PSL}_2(\Z)$.

Karpenkov's Problem 19 asks for a \emph{natural generalization of this tessellation to higher-dimensional hyperbolic geometry}, that is, constructions in $\mathbb H^d$ for $d\geq3$ together with an explanation of how they generalize the Farey construction. This is a construction and comparison problem rather than a fully quantified conjecture. The source does not specify which features are required: ideal vertices, a rational or arithmetic boundary structure, a determinant-type incidence criterion, a symmetry group, a tessellation by simplices versus other cells, or compatibility across dimensions. These possibilities identify missing choices; none is imposed here as an additional source hypothesis.

The supplied research summary reports partial progress through Farey-type tessellations of $\mathbb H^3$, arithmetic hyperbolic tessellations, and relations between Klein sails and cone-Hilbert geometry. It says no canonical universally agreed generalization in all dimensions has been established. No explicit construction, reference, or uniqueness criterion is supplied for these claims. Consequently ``natural'' and ``canonical'' cannot presently be converted into a single existence or uniqueness assertion from this record alone.
\subsection{Short English statement}
The task is to extend the arithmetic ideal-triangle Farey tessellation of the hyperbolic plane to higher-dimensional hyperbolic spaces. Several related constructions are reported, but the source does not specify which Farey features a satisfactory generalization must preserve.
\subsection{Sources}
\begin{itemize}
\item Local UnsolvedMath registry, permanent ID 1900019, exact source code catalogue entry, ``Geometry of Continued Fractions --- Further open questions''; status partially solved, classification PARTIAL-PROGRESS.
\item Karpenkov, \emph{Open problems in geometry of continued fractions} (2017), Problem 19, \url{https://arxiv.org/abs/1712.01450}. The record identifies extraction from public arXiv source TeX, accessed July 29, 2026, and says the source presented the item as open.
\item The supplied research summary mentions Farey-type tessellations of $\mathbb H^3$, arithmetic hyperbolic tessellations, and Klein-sail/cone-Hilbert connections as partial responses. It gives no named author, explicit construction, or authoritative citation for these claims, and its assessment about a universally agreed all-dimensional construction is not a precise theorem.
\end{itemize}
\end{document}
